\section{\texorpdfstring{Ramification of $G$-Torsors}{Ramification of G-Torsors}}\label{section:ramification_of_torsors}

Generally speaking, assume we are given a locally Noetherian scheme $X$,
a dense open $U \subset X$, a finite \'etale morphism $f: Y \to U$ and a prime divisor $Z \subset X$ such that $Z \cap U = \emptyset$
and the local ring $\Oo_{X, \xi}$ of $X$ at the generic point $\xi$ of $Z$ is a discrete valuation ring.
Setting $K_\xi = \Frac(\Oo_{X, \xi})$ we obtain a cartesian square
\[ \xymatrix{ \mathop{\mathrm{Spec}}(K_\xi ) \ar[r] \ar[d] & U \ar[d] \\ \mathop{\mathrm{Spec}}(\mathcal{O}_{X, \xi }) \ar[r] & X } \]
of schemes.
In particular, we see that $Y \times _ U \mathop{\mathrm{Spec}}(K_\xi )$ is the spectrum of a finite separable algebra $L_\xi /K_\xi $.
\begin{definition}\label{definition:tamely_ramified_in_codim_1}
We say:
\begin{enumerate}
    \item $Y$ is \textit{unramified} over $X$ at $(Z, \xi)$ resp. $Y$ is \textit{tamely ramified} over $X$ at $(Z, \xi)$
        if $L_\xi /K_\xi $ is unramified, resp. tamely ramified with respect to $\mathcal{O}_{X, \xi }$
        (\Cref{definition:tamely_ramified_dvrs})

    \item $Y$ is unramified over $X$ in codimension $1$, resp. $Y$ is tamely ramified over $X$ in codimension $1$ if $L_\xi /K_\xi $ is unramified, resp. tamely ramified
    with respect to $\mathcal{O}_{X, \xi }$ for every $(Z, \xi )$ as above.
\end{enumerate}

More precisely, we decompose $L_\xi $ into a product of finite separable field extensions of $K_\xi $ and we require each of these to be unramified, resp. tamely ramified with
respect to $\mathcal{O}_{X, \xi }$.

\end{definition}

Now back to our original assumptions, $C$ is a projective smooth geometrically connected curve
over a field $k$ and $U=C\setminus \mdls$. Every $(Z, \xi)$ is now a closed point $P$ and $\Oo_{C, P}$ is a discrete valuation ring.
Assume $f: Y \to U$ is actually a $G$-torsor $\cP \to U$ in $U_{\text{\'Et}}$ for $G$ a finite abelian group.

By passing to the completion $\widehat{\Oo}_{C,P}$, we obtain the complete local field $\widehat{K}_P$.
Restricting the torsor $\cP$ to the punctured formal neighborhood $\Spec(\widehat{K}_P)$
 allows us to study the ramification in a strictly local context.
 This restriction yields a $G$-torsor over a field, which necessarily decomposes into a disjoint union of spectra of finite separable field extensions
 $$ \cP|_{\Spec(\widehat{K}_P)} \cong \bigsqcup \Spec(F) $$
 such that:
 \begin{enumerate}
    \item These extensions $F/\widehat{K}_P$ are all isomorphic and Galois.
    \item The Galois group $H = \Gal(F/\widehat{K}_P)$ identifies with the stabilizer of any connected component of the pullback $\cP|_{\Spec(\widehat{K}_P)}$,
     which is a subgroup of $G$.
    \item  The number of such components is then given by the index $[G:H]$.
 \end{enumerate}

 \begin{definition}\label{definition:local_ramification_boundness}
    We say that the extension $F/\widehat{K}_P$ \textit{has ramification bounded by $r$} if the ramification group $H^r$ is trivial.
    We say \textit{$\cP$ has ramification bounded by $r$ at $P$} if the corresponding local field extensions $F/\widehat{K}_P$
    satisfy this condition.
 \end{definition}


\begin{definition}
    Let $\mdls = \sum n_i P_i$ be a modulus on $C$. We say the $G$-torsor $\cP \to U$ has
    \textbf{ramification bounded by $\mdls$} if, for each point $P_i$, the local restriction
    $\cP|_{\Spec(\widehat{K}_{P_i})}$ has ramification bounded by the coefficient $n_i$.
\end{definition}

This local property remains invariant under base change to the separable closure.
Let $\bar{s} = \Spec(\bar{k}) \to \Spec(k)$
be a geometric point.
The higher ramification groups for $\cP|_{\Spec(\widehat{K}_P)}$ are isomorphic to those of
the geometric restriction $\cP_{\bar{k}}$ at any point $\bar{x}$ lying over $P$.
Thus an equivalent definition:
\begin{definition}
    Let $\mdls = \sum n_i P_i$ be a modulus on $C$.
    The $G$-torsor $\cP \to U$ has \textbf{ramification bounded by $\mdls$}
    if for every geometric point $\bar{x}$ of $\mdls$ (with image $\bar{s}$ in $\Spec(k)$),
    the restriction of $\cP$ to$$ \Spec(\widehat{\mathcal{O}}_{C_{\bar{k}},\bar{x}}) \times_{C_{\bar{k}}} U_{\bar{k}} $$
    has ramification bounded by the multiplicity of $\mdls_{\bar{s}}$ at $\bar{x}$ in the sense of
    \Cref{definition:local_ramification_boundness}.
\end{definition}
These definitions are equivalent because $\widehat{\mathcal{O}}_{C_{\bar{k}},\bar{x}} \cong \widehat{\mathcal{O}}_{C,P} \otimes_k \bar{k}$. Furthermore, the notions of tame ramification and unramifiedness derived here coincide with the codimension-$1$ criteria in \Cref{definition:tamely_ramified_in_codim_1}.

There is another equivalent formulation through the language of characters.
The group $G$, being a finite abelian group over $k$, is \'etale.
Consequently, the correspondence between $G$-torsors and the fundamental group allows us to describe $G$-torsors
 via characters. Specifically, there is an isomorphism (\Cref{cor:abelian_equivilance_fundamental_torsors}):
\[ \mathrm{Hom}_{\mathrm{cont.}}(\pi_1^{\text{\'et}}(U, \overline{x}), G) \xrightarrow{\cong} \mathrm{Tors}(U_{\text{\'et}}, G) \]

In our local setting, where we consider the restriction to the complete local field $\widehat{K}_P$,
the fundamental group identifies with the absolute Galois group
$G_{\widehat{K}_P} := \mathrm{Gal}({\widehat{K}_P}^{\mathrm{sep}}/{\widehat{K}_P})$.
Thus, the restricted $G$-torsor $\mathcal{P}|_{\mathrm{Spec}(\widehat{K}_P)}$ corresponds to a unique continuous homomorphism:
\[ \rho : G_{\widehat{K}_P} \to G \]

Under this correspondence, the torsor $\mathcal{P}$ has ramification bounded by $r$ at $P$
if and only if the homomorphism $\rho$ trivializes the $r$-th ramification group in the upper numbering:
\[ \rho(G_{\widehat{K}_P}^r) = \{ 1 \} \]

\subsection*{Basic Properties of Ramification of \texorpdfstring{$G$-Torsors}{G-Torsors}}
We now prove some elementary lemmas regarding the ramification of $G$-torsors.
\begin{lemma}\label{lemma:bounding_ramification_of_contracted_product}
    Let $G$ be a finite abelian group and $X$ be a locally Noetherian scheme over a field $k$.
    Let $U \subset X$ be a dense open subset and let $(Z, \xi)$ be a prime divisor in the complement $X \setminus U$.
    Assume $\cP_1$ and $\cP_2$ are two $G$-torsors on $U_{\text{ét}}$, such that $\cP_1$ has ramification bounded by $r_1$ at $(Z, \xi)$, and $\cP_2$ has ramification bounded by $r_2$ at $(Z, \xi)$.
    Then their product $\cP_1 \otimes^G \cP_2$ has ramification bounded by $\max(r_1, r_2)$ at $(Z, \xi)$.
\end{lemma}
\begin{proof}
    Let $A=\widehat{\Oo_{X, \xi}}$ and let $K=\Frac(A)$.
    Let $\rho_1, \rho_2: G_K \to G$ be the associated continuous homomorphisms corresponding to the $G$-torsors
    $\cP_1|_{\Spec(K)}$, $\cP_2|_{\Spec(K)}$.
    We finish by noticing that the associated character of $(\cP_1 \otimes^G \cP_2)|_{\Spec(K)}$ is $\rho=\rho_1 + \rho_2$.
\end{proof}

\begin{lemma}\label{lemma:descend_ramification_along_etale}

Let $f: X \to Y$ be a morphism of locally Noetherian schemes. Let $U_X \subset X$ and $U_Y \subset Y$ be dense open subschemes such that $f^{-1}(U_Y) \subset U_X$.

Let $(Z_X, \xi_X)$ and $(Z_Y, \xi_Y)$ be prime divisors of $X$ and $Y$ such that $Z_X \cap U_X = \emptyset$ and $Z_Y \cap U_Y = \emptyset$.
Suppose $f(\xi_X) = \xi_Y$ and that $f$ is étale at $\xi_X$.
Let $\mathcal{P}$ be a $G$-torsor over $U_Y$, and let $f^{-1} \mathcal{P}$ be its pullback to $f^{-1}(U_Y) \subset U_X$.
Then, the ramification of $f^{-1} \mathcal{P}$ at $\xi_X$ is bounded by $r$ if and only if the ramification of $\mathcal{P}$ at $\xi_Y$ is bounded by $r$.
\end{lemma}
\begin{proof}
The boundedness of ramification is determined by the behavior of the torsor over the completion of the local rings at the generic points.
Let $A = \mathcal{O}_{Y, \xi}$ and $B = \mathcal{O}_{X, \xi_X}$ be the discrete valuation rings at the generic points, with fraction fields
$K$ and $L$ respectively. Since $f$ is étale at $\xi_X$, the map $A \to B$ is a flat, unramified (hence weakly unramified) local homomorphism.
Consequently, the extension of completions $\widehat{L} / \widehat{K}$ is a finite unramified extension of complete discretely valued fields.
We finish by noticing that the upper numbering filtration on the absolute Galois group is compatible with unramified base change.\footnote{
 Specifically, let $G_K = \operatorname{Gal}(K^{sep}/K)$ and $G_L = \operatorname{Gal}(L^{sep}/L)$. For an unramified extension, the Herbrand function is the identity, which implies that for any $r \geq 0$:
\[ G_L^r = G_K^r \cap G_L \]}
\end{proof}

% --- NEW (not in the current chapter): two further stability lemmas, needed
%     for the assembly of the general abelian case in 2.5.4. Their proofs are
%     character computations in the style of the two lemmas above. ---

\begin{lemma}\label{lemma:ramification_of_fiber_product}
    Let $G_1, G_2$ be finite abelian groups and let $X, U, (Z, \xi)$ be as in
    \Cref{lemma:bounding_ramification_of_contracted_product}. Let $\cP_1$ be a
    $G_1$-torsor and $\cP_2$ a $G_2$-torsor on $U_{\text{ét}}$, and consider the fiber
    product $\cP_1 \times_U \cP_2$, which is a $(G_1 \times G_2)$-torsor on
    $U_{\text{ét}}$. If $\cP_i$ has ramification bounded by $r_i$ at $(Z, \xi)$ for
    $i = 1, 2$, then $\cP_1 \times_U \cP_2$ has ramification bounded by
    $\max(r_1, r_2)$ at $(Z, \xi)$. Moreover, $\cP_1 \times_U \cP_2$ is tamely
    ramified (resp.\ unramified) at $(Z, \xi)$ if and only if both $\cP_1$ and
    $\cP_2$ are.
\end{lemma}
\begin{proof}
    Let $A=\widehat{\Oo_{X, \xi}}$, $K=\Frac(A)$, and let
    $\rho_i: G_K \to G_i$ be the characters of $\cP_i|_{\Spec(K)}$. The
    character of $(\cP_1 \times_U \cP_2)|_{\Spec(K)}$ is
    $\rho = (\rho_1, \rho_2) : G_K \to G_1 \times G_2$, since torsors under a
    product group correspond to pairs of characters,
    $H^1(K, G_1 \times G_2) = H^1(K, G_1) \times H^1(K, G_2)$, compatibly with
    the projections. Now for any closed subgroup $N \subset G_K$ (we take the
    upper-numbering group $G_K^r$, resp.\ the wild inertia subgroup, resp.\ the
    inertia subgroup) we have $\rho(N) = \{1\}$ if and only if
    $\rho_1(N) = \{1\}$ and $\rho_2(N) = \{1\}$.
\end{proof}

\begin{lemma}\label{lemma:ramification_of_induced_torsor}
    Let $\iota : H \hookrightarrow G$ be an injective homomorphism of finite abelian
    groups, let $X, U, (Z, \xi)$ be as above, let $\cQ$ be an $H$-torsor on
    $U_{\text{ét}}$, and let $\iota_* \cQ := (\cQ \times G)/H$ be the induced
    $G$-torsor (where $H$ acts by $h \cdot (q, g) = (qh,\ \iota(h)^{-1}g)$).
    Then $\iota_* \cQ$ has ramification bounded by $r$ at $(Z, \xi)$ if and only if
    $\cQ$ does; likewise for tame ramification and for unramifiedness.
\end{lemma}
\begin{proof}
    Let $A=\widehat{\Oo_{X, \xi}}$, $K=\Frac(A)$, and let $\chi: G_K \to H$ be the
    character of $\cQ|_{\Spec(K)}$; the character of $(\iota_*\cQ)|_{\Spec(K)}$ is
    $\iota \circ \chi$. Since $\iota$ is injective, $\iota \circ \chi$ and $\chi$
    have the same kernel, so for any closed subgroup $N \subset G_K$ we have
    $(\iota \circ \chi)(N) = \{1\}$ if and only if $\chi(N) = \{1\}$.
    (Concretely: the connected components of the two torsors over $\Spec(K)$ are
    spectra of the \emph{same} field extension $F/K$, the fixed field of
    $\ker \chi$; the induced torsor merely has $[G : H]$ times as many components.)
\end{proof}
